\chapter{Testing and Quality Assurance}

\section{Test strategy}
\begin{longtable}{L{2.6cm}L{6.6cm}L{4.6cm}}
\caption{Test levels.}\label{tab:testlevels}\\
\toprule
Level & What is tested & How \\
\midrule\endfirsthead
\toprule Level & What is tested & How \\ \midrule\endhead
Unit & parser on a real recording and on garbage; safe unpickler against a malicious pickle; every noise function (shape, determinism, no in-place change, effect size); every model's forward pass with empty segments & \texttt{pytest tests/} \\
Data validation & split disjointness (RadHAR recordings, MiliPoint subjects); tensor shapes and finiteness; deterministic evaluation samples; train-only statistics & \texttt{pytest} + leakage audit \\
Integration & exported bundle $\rightarrow$ raw recording $\rightarrow$ probabilities (sum to 1, deterministic); short input rejected & \texttt{pytest}, CLI \\
Checkpoint & bundle reload without retraining; seed selected by validation & CLI + demo \\
Application & app start, sample loading, preprocessing view, inference, all 7 noise conditions, sample switching, error handling & Streamlit \texttt{AppTest}, headless-browser screenshots \\
Performance & latency (median of 200, batch 1, 1 thread), FLOPs, model size & \texttt{train.py} \\
Reproducibility & identical robustness results on repeated evaluation & robustness determinism check \\
\bottomrule
\end{longtable}

\section{Results}
\begin{table}[h]
\centering
\caption{Test results (generated by \texttt{scripts/make\_qa.py}).}
\label{tab:testresults}
\begin{tabular}{lrr}
\toprule
Suite & Passed & Failed \\
\midrule
Unit / data / integration (\texttt{pytest tests/}) & 23 & 0 \\
Demo application (\texttt{scripts/apptest\_demo.py}, 28 interaction steps) & 28 & 0 \\
Leakage audit checks (\texttt{scripts/leakage\_audit.py}) & 9 & 0 \\
\bottomrule
\end{tabular}
\end{table}
The demo test drives the real application headlessly: it loads both exported models, every bundled test sample and every perturbation condition, and fails on any exception or error message shown to the user.

\section{Quality-assurance matrix}
\begin{longtable}{L{2.5cm}L{5.6cm}L{4.4cm}L{1.4cm}}
\caption{Final quality-assurance matrix (evidence-based).}\label{tab:qa}\\
\toprule
Area & Check & Evidence & Status \\
\midrule\endfirsthead
\toprule Area & Check & Evidence & Status \\ \midrule\endhead
\csname @@input\endcsname ../paper/generated/qa_rows.tex 
\bottomrule
\end{longtable}
